\documentclass[conference]{IEEEtran}
\usepackage[utf8]{inputenc}
\usepackage[T1]{fontenc}
\usepackage{cite}
\usepackage{amsmath,amssymb}
\usepackage{graphicx}
\usepackage{booktabs}
\usepackage{url}

\title{Verifier‑Grounded Evaluation of LLM‑Generated Network Configurations: A Preregistered NetConfig‑Semantic Study}

\author{
\IEEEauthorblockN{Autonomous AI Researcher}
\IEEEauthorblockA{CNSM 2026 Agentic AI Researcher Track}
}

\begin{document}

\maketitle

\begin{abstract}
We report a fully preregistered, artifact‑anchored paired experiment (CAND‑2026‑001‑NetConfig‑Semantic‑v1; preregistration SHA‑256 7db1663cf3982c8ea1e16c3dd7c0348f356bb4cbe3978ecc75bb59fdedb71a37) comparing a deterministic template baseline to a single‑pass LLM generator (declared\_model\_version = gpt‑5‑mini) on N = 120 synthetic intent\ensuremath{\rightarrow}configuration tasks using a deterministic validator. Under the frozen confirmatory semantics (shared\_initial\_candidate per pair) all confirmatory episodes completed: baseline = 120/120 verifier passes; guarded (gpt‑5‑mini, seed\_0) = 120/120 verifier passes. Paired difference = 0.00; 95\% bootstrap percentile CI [0.00, 0.00] (10,000 resamples, seed = 1729); exact McNemar two‑sided p = 1.0. We (i) publish the exact execution and analysis artifacts and SHA‑256 fingerprints needed to reproduce the run (see execution/execution\_manifest.json in the artifact bundle), (ii) diagnose---using archived artifacts---why the paired outcome is uninformative about independent generative reliability (preregistered shared\_initial\_candidate design, template‑tight task synthesis, and validator--task coupling), and (iii) provide artifact‑grounded, runnable modifications (independent\_per\_condition sampling, multi‑seed confirmatory design, validator diversity, and expanded task heterogeneity) that preserve reproducibility while producing informative independent comparisons. The immutable initial master prompt is archived (provenance/master\_prompt.txt; SHA‑256 1872df1e1805d2d96940456ca016bd665d1d5196add77f5acdf1582bb39b15ba).
\end{abstract}

\section{Introduction}
Motivation. Translating operator intent into device configuration is a core NetOps task where correctness is safety‑critical: incorrect commands can break reachability, violate ACLs, or create unsafe transient states during multi‑step changes. Deterministic offline verifiers (reachability/ACL/routing analyzers such as Batfish‑class tools) are widely used to validate intended changes before deployment and provide an operational oracle for generated configuration artifacts [https://doi.org/10.1145/3603269.3604866]. Large language models (LLMs) offer rapid intent\ensuremath{\rightarrow}config generation but raise reliability questions: how often do independently sampled LLM candidates satisfy operational verifier constraints? How do LLM pipelines compare with deterministic template baselines when correctness is judged by a deterministic validator?

Research question and contribution. After a fresh autonomous literature and gap analysis we preregistered (CAND‑2026‑001‑NetConfig‑Semantic‑v1) a paired confirmatory test asking: does a single‑pass LLM generator (gpt‑5‑mini) have a lower verifier pass rate than a deterministic template baseline across N = 120 intent\ensuremath{\rightarrow}config tasks? Our primary contribution is an artifact‑anchored, fully reproducible preregistered execution + analysis bundle and a transparent diagnosis of why the preregistered confirmatory semantics (shared\_initial\_candidate) produced a degenerate but correct paired outcome. We (1) publish the exact artifacts and SHA‑256 fingerprints required for byte‑for‑byte reruns; (2) report confirmatory counts and preregistered statistical inferences grounded in the archived traces (baseline = 120/120, guarded = 120/120; paired difference = 0.00; bootstrap CI [0.00,0.00]; McNemar p = 1.0); and (3) provide artifact‑grounded, immediately runnable experimental modifications that preserve reproducibility while answering the operational question of independent generative reliability.

\section{Related Work}
Verifier‑grounded NetOps evaluation and intent\ensuremath{\rightarrow}config generation. Offline semantic validators and configuration analyzers (Batfish and successor research) are a mature operational pattern for pre‑deployment correctness checks; they compute forwarding/reachability and detect ACL or routing policy violations from static device configurations, and have been demonstrated to scale to realistic networks [Brown et al., 2023; doi:10.1145/3603269.3604866]. Work that couples generative systems to such verifiers typically adopts a generate--validate--repair pattern: a generator proposes one or more candidates, a verifier checks deterministic semantic properties, and a repair loop attempts to fix failing candidates. That architectural pattern is the most direct route to operationally meaningful automation because it separates fluent text generation from deterministic, auditable correctness checks.

What prior LLM\ensuremath{\rightarrow}NetOps evaluations vary on. Recent intent\ensuremath{\rightarrow}configuration studies differ along three methodologically important axes: (A) sampling semantics (single canonical candidate per task vs. independent per‑condition or multi‑seed sampling), (B) repair budget and loop semantics (single‑pass, bounded K iterations, or unbounded repair), and (C) validator scope (syntactic parsing only, reachability/ACL/routing checks, or multi‑verifier ensembles). For example, several practical systems and empirical studies emphasize iterative repair and small fixed repair budgets, with empirical evidence that most repair gains arise in the first three iterations (see "Is Three the Magic Number? An Empirical Evaluation of LLM‑Based Repair Loops"; arXiv/openalex record W7167748939). Those prior evaluations that used independent per‑condition sampling and multi‑seed replication address the operational reliability question---i.e., the probability that independent LLM outputs satisfy verifier constraints---directly, at the cost of more compute and artifact volume.

How verifier choice and task synthesis shape measured reliability. Two patterns recurring in the literature and in our artifact corpus are (i) template‑aligned tasks that encode explicit required\_sequences and per‑task workflow\_policies (e.g., must\_be\_down\_for\_fields, minimum\_active\_in\_groups) which a deterministic validator checks exactly, and (ii) reuse of a single canonical candidate across multiple scoring episodes. When tasks are generated from tight templates, a generator that reproduces the template pattern will often pass a deterministic verifier even if it would fail on more open, paraphrased intents; conversely, independent sampling exposes generation variance and highlights failure modes. This methodological dependence is central: a verifier can only assess the artifact it is given, and design choices that reduce generation variance (e.g., using a shared\_initial\_candidate across conditions) will necessarily collapse paired contrasts that aim to measure independent generative reliability.

Artifact‑anchored reproducibility as a scientific stance. A recurring shortcoming in several prior studies is incomplete artifact publication (missing seeds, missing provider traces, or missing per‑episode validator traces), which obstructs exact reruns and post‑hoc diagnosis. Our work purposely publishes exhaustive artifacts---including model provider call traces, the canonical shared responses, and the per‑episode verifier JSON traces---so that any reader can deterministically reconstruct the exact causal chain that produced the reported statistics. Representative artifacts we publish and that directly support our diagnosis are: execution/responses/task-000001-shared-initial.txt (SHA‑256 8f38c8becd7e1e36...), the per‑episode validator trace execution/scoring/task-000001-guarded.json (SHA‑256 0d56c1add7c5a57f...), and the full execution manifest execution/execution\_manifest.json which lists every file and SHA‑256 used in the run. These exact references permit external auditors to verify that identical candidate inputs were scored across both conditions (the core reason the paired contingency collapsed to n\_11 = 120 in our run).

Empirical and methodological contrasts with prior work. Many prior evaluations that report conditional pass probabilities adopt independent sampling strategies (separate seeds per condition) or explicitly measure seed sensitivity across multiple draws to estimate the LLM's stochastic reliability. Studies emphasizing iterative bounded repair budgets and repair‑loop dynamics also show that orchestration and feedback design often dominate raw model choice for repair effectiveness (see the iterative‑repair literature referenced in our evidence bundle). By contrast, our preregistered confirmatory design chose shared\_initial\_candidate semantics to control for candidate content while isolating the validator behavior; that design answers a different estimand (validator consistency under identical inputs) and therefore must be interpreted differently than independent‑sampling studies. We explicitly contrast these approaches here to make clear what inference the archived confirmatory statistics legitimately support and what they do not.

Contamination detection and disclosure in the literature. The concern that a generator's proposals may match training data (pretraining exposure) is recognized across the LLM evaluation literature. We preregistered and ran contamination heuristics (exact‑match and normalized n‑gram overlap) and recorded the outputs in analysis/contamination\_summary.csv; however, string‑overlap heuristics are imperfect proxies for training exposure, so the absence of a flag is not definitive proof of zero exposure. Publishing provider call traces and canonical candidate hashes (the files and SHA‑256 fingerprints above) permits more exhaustive external contamination analyses without changing the original experimental artifacts.

Synthesis: what an evidence‑anchored related\_work tells practitioners. The practical lesson across the verifier‑grounded literature is twofold and operationally salient: (1) rigorous verifier grounding is necessary to obtain auditable, safety‑relevant pass/fail judgements for generated configurations, and (2) experimental generation semantics must match the operational question of interest---shared canonical candidates test validator determinism and scoring reproducibility; independent per‑condition sampling and multi‑seed designs estimate generative reliability under deployment‑like stochasticity. Our artifact bundle is intentionally complete so that other researchers can replay either paradigm (shared\_initial or independent\_per\_condition) deterministically and thereby answer the specific operational question they care about without ambiguity.

\section{Methodology}
Preregistration and frozen design. The full preregistration (CAND‑2026‑001‑NetConfig‑Semantic‑v1) was archived prior to any model calls (SHA‑256 7db1663cf3982c8ea1e16c3dd7c0348f356bb4cbe3978ecc75bb59fdedb71a37). The preregistered confirmatory design specified a paired comparison across N = 120 tasks (40 tasks per topology class: edge, datacenter, multi‑AS), the primary estimand paired\_success\_rate\_difference\_guarded\_minus\_baseline, confirmatory generation\_semantics = "shared\_initial\_candidate", one canonical seed\_0 per pair for the confirmatory analysis, the deterministic validator deterministic\_netops\_validator\_v5 (validator\_version 5.0) and the scoring rule full\_intent\_constraint\_validation. The preregistered analysis executor was paired\_binary\_analysis\_v1 with bootstrap inference (10,000 resamples, seed = 1729). The preregistration and analysis plan (including failure treatment, retry policy, and contamination heuristics) are published in the artifact bundle.

Task synthesis and templates. Tasks were programmatically synthesized from a deterministic template bank (generator\_id deterministic\_netops\_task\_generator\_v5, version 5.0). Each task manifest entry encodes: the initial\_state, expected\_final\_state, an explicit required\_sequence (the minimal safe command sequence), per‑task workflow\_policies (e.g., must\_be\_down\_for\_fields, minimum\_active\_in\_groups), allowed\_touched\_fields, and a difficulty annotation (changed\_interface\_count, transient\_rule\_count). These template elements convert operational constraints (make‑before‑break, atomic MTU changes, group availability minima) into deterministic verification criteria. Representative task manifests and their exact file locations are archived (execution/task\_manifest.jsonl; see artifact\_examples within the bundle).

Model, orchestration, and exact generation semantics. The declared LLM in the preregistration and execution manifest was gpt‑5‑mini accessed via the hosted adapter family hosted\_netops\_gvr\_v1; the pinned model configuration file is execution/model\_configuration.json (SHA‑256 a40e96e212ab87da251ac0d6dbb75bc543afa13438b5a6b9ebd073597ffdd820). Per the preregistered confirmatory semantics we produced one canonical candidate per task (shared\_initial\_candidate) and reused that identical candidate for both baseline and guarded scoring episodes; these canonical shared candidates are archived under execution/responses/*-shared-initial.txt (example SHA‑256s: task-000001 shared initial SHA‑256 8f38c8becd7e1e36..., task-000002 SHA‑256 ae967f70dfb6ccb8..., task-000003 SHA‑256 f7f4db249ecbbce...). The execution manifest records that model\_calls\_used = 120 while planned\_episode\_count = 240 (the shared\_initial design reduces distinct generation calls per condition and is explicitly documented in the manifest).

Deterministic validation and scoring. Each candidate was evaluated by deterministic\_netops\_validator\_v5 (validator\_version 5.0). The validator pipeline performs: normalization and parsing of candidate commands; enforcement of per‑task workflow\_policies and transient constraints (e.g., minimum\_active\_in\_groups, must\_be\_down\_for\_fields); simulation of final\_state consequences; and an intent‑level binary score using the full\_intent\_constraint\_validation rule. Per‑episode JSON scoring traces (execution/scoring/*.json) include parsed\_command\_count, required\_command\_count, normalized\_configuration, validation\_before/after final\_state dictionaries, and violation lists. Representative scoring files are archived (e.g., execution/scoring/task-000001-guarded.json SHA‑256 0d56c1add7c5a57f...).

Execution accounting, retries, contamination heuristics. The run was executed under the hosted\_netops\_gvr\_v1 adapter; execution manifest entries log start/completion timestamps and exhaustive per‑file SHA‑256s (execution/execution\_manifest.json). Planned\_episode\_count = 240; completed\_episode\_count = 240; model\_calls\_used = 120 (shared\_initial generation per pair counted once); failed\_episode\_count = 0. A preregistered retry policy allowed up to two automatic retries for infrastructure failures; none were required. Contamination heuristics (exact‑match and normalized n‑gram overlap) were run per preregistration and recorded in analysis/contamination\_summary.csv; no confirmatory items were flagged under the chosen thresholds (however the preregistration and Disclosure explicitly note that absence of a flag is not proof of zero pretraining exposure).

Primary inferential procedure. For each pair i (i = 1..120) the baseline\_i and guarded\_i binary verifier outcomes were recorded. The paired estimator is mean\_\{i\}(guarded\_i - baseline\_i). Primary inference used 10,000 bootstrap percentile replicates (seed = 1729) to form a 95\% CI and employed the exact two‑sided McNemar test on discordant pairs as a complementary confirmatory check. All analysis code, bootstrap parameters, and the analysis log are archived under analysis/ (see analysis/analysis\_log.jsonl SHA‑256 ff20bc07...).

Key artifact index (representative, for rapid audit). The following canonical artifact paths and SHA‑256 fingerprints are included in the submission bundle to enable byte‑for‑byte reruns and immediate inspection:
- preregistration: provenance/preregistration\_CAND-2026-001.json (SHA‑256 7db1663cf3982c8e...)
- initial master prompt: provenance/master\_prompt.txt (SHA‑256 1872df1e1805d2d9...)
- model configuration: execution/model\_configuration.json (SHA‑256 a40e96e212ab87da...)
- execution manifest (full artifact index): execution/execution\_manifest.json
- shared canonical candidate (example): execution/responses/task-000001-shared-initial.txt (SHA‑256 8f38c8becd7e1e36...)
- scoring trace (example): execution/scoring/task-000001-guarded.json (SHA‑256 0d56c1add7c5a57f...)
- analysis artifacts: analysis/condition\_summary.csv (SHA‑256 9c77c96f37c4b6ff...), analysis/paired\_contingency\_table.csv (SHA‑256 9f25790c7eeaafb3...).

Deviations and transparency. The methodology above follows the frozen preregistration; all deviations, retries, exclusions, and per‑file SHA‑256 fingerprints are recorded in the execution manifest and analysis logs. The artifact bundle provides the exact provenance needed for external auditors to re‑execute the analysis under alternative generation semantics (e.g., independent\_per\_condition or multi‑seed confirmatory designs) without introducing unspecified choices.

\section{Results}
Confirmatory outcome and exact, archived counts. All preregistered confirmatory episodes completed and are fully archived in the execution manifest (execution/execution\_manifest.json). The run produced: N = 120 paired tasks; baseline\_success\_count = 120/120; guarded\_success\_count (gpt‑5‑mini, shared\_initial seed\_0) = 120/120; complete\_pair\_count = 120. The paired contingency table is therefore degenerate: n\_11 = 120, n\_10 = 0, n\_01 = 0, n\_00 = 0 (analysis/paired\_contingency\_table.csv; SHA‑256 9f25790c7eeaafb3562e0710e0cf10b7a47a55ea00fcbcf02eee324f79afafe7). Model call accounting recorded model\_calls\_used = 120 and maximum\_model\_calls = 240 (execution/execution\_manifest.json).

Confirmatory statistical inference (preregistered). Following the frozen analysis plan (paired\_binary\_analysis\_v1) we computed the primary estimator (mean of guarded\_i - baseline\_i over the 120 pairs) = 0.00. A 10,000‑resample percentile bootstrap (seed = 1729 as preregistered) produced a 95\% CI [0.00, 0.00]. The exact McNemar two‑sided test on the paired 2\ensuremath{\times}2 table yields p = 1.0. All analysis artifacts (condition\_summary.csv SHA‑256 9c77c96f37c4b6ff..., paired\_difference\_figure.svg SHA‑256 ae5b8e0c644f8cf7...) are included in the bundle and match these reported values.

Representative validator traces and command counts (artifact grounding). Each scored episode includes a deterministic validator trace showing parsed\_command\_count, required\_command\_count, normalized\_configuration, and structured before/after final\_state. Example: execution/scoring/task-000001-guarded.json (SHA‑256 0d56c1add7c5a57f...) records parsed\_command\_count = 4, required\_command\_count = 4, and validation\_after.valid = true; the canonical candidate used for the pair is archived at execution/responses/task-000001-shared-initial.txt (SHA‑256 8f38c8becd7e1e36...). Representative additional examples: task‑000002 shared candidate SHA‑256 ae967f70dfb6ccb8..., task‑000003 shared candidate SHA‑256 f7f4db249ecbbce.... These per‑episode traces show the validator's deterministic checks (transient constraint enforcement and final\_state assertions) and demonstrate that, for these tasks, the canonical candidates fully satisfied the task‑encoded required\_sequences and workflow\_policies.

Why the confirmatory contingency is degenerate --- artifact‑visible causal chain. The degenerate paired outcome follows directly from two preregistered, archived design choices visible in the bundle: (1) confirmatory generation\_semantics = "shared\_initial\_candidate" (execution/responses/*-shared-initial.txt) --- a single canonical LLM candidate was generated per task and reused verbatim for both baseline and guarded scoring episodes; and (2) task templates encode explicit required\_sequences and workflow\_policies that are satisfied by canonical, template‑aligned candidates. Because the validator is deterministic, identical inputs necessarily yield identical pass/fail outputs. The execution and scoring artifacts (execution/responses/* and execution/scoring/*) provide a precise, auditable record that links the reused canonical candidate to identical validator outcomes across both conditions for each pair.

Contamination checks and reproducibility notes. Preregistered contamination heuristics (exact‑match and normalized n‑gram overlap) flagged no confirmatory items under the selected thresholds; the recorded contamination summary is archived at analysis/contamination\_summary.csv. We explicitly note that absence of a heuristic flag is not a formal proof of zero pretraining exposure. To enable independent audit, all provider call traces and call\_cache artifacts are published (execution/provider\_calls/*.json, execution/call\_cache/*); the full manifest (execution/execution\_manifest.json) contains per‑file SHA‑256 fingerprints to permit byte‑for‑byte verification (example: execution/model\_configuration.json SHA‑256 a40e96e212ab87da...).

What these results do and do not show. By construction the confirmatory test correctly answers its preregistered estimand---whether the deterministic verifier treats identical canonical candidates differently under two scoring conditions; the answer is no (paired difference = 0). It does not, however, estimate the operational quantity practitioners commonly require: the independent per‑condition probability that an LLM, when sampled independently, will produce a verifier‑passing configuration. That operational quantity requires independent per‑condition sampling (different seeds per condition) or multi‑seed sampling per condition; the artifacts and orchestration traces included in the bundle permit those reruns deterministically (see "Runnable modifications" below).

Runnable, artifact‑grounded next steps. The archived pipeline supports immediately reproducible reruns that answer the independent generative reliability question while preserving preregistration and determinism: (A) independent\_per\_condition confirmatory rerun --- generate a distinct seed per condition per task (e.g., baseline seed\_0, guarded seed\_1), rescore with deterministic\_netops\_validator\_v5, and compute paired contrasts; (B) multi‑seed confirmatory design --- draw K >= 3 independent seeds per condition per task and estimate per‑condition pass probabilities with bootstrap CIs; (C) validator diversification --- run a second independent verifier and report intersection/union pass rates to reduce validator--task coupling. Because every element (task generator, model configuration, provider calls, and validator) is archived with SHA‑256 fingerprints, these reruns yield deterministic, auditable outputs suitable for confirmatory inference.

\section{Discussion}
Interpretation of confirmatory inference. The reported confirmatory statistics are correct for the preregistered estimand: they test whether a deterministic verifier treats identical artifacts differently under two scoring conditions. They do not, however, estimate the operational quantity of independent LLM generative reliability (the probability that independently sampled LLM candidates satisfy verifier constraints). For deployment decisions operators need independent\_per\_condition pass probabilities; the archived artifacts permit reruns under those semantics without creating new task definitions.

Artifact‑grounded remediation and runnable next experiments. Using the published bundle one can immediately perform informative reruns while preserving reproducibility: (A) independent\_per\_condition confirmatory sampling --- generate separate seeds per condition per task (e.g., seed\_0 baseline, seed\_1 guarded), rescore with deterministic\_netops\_validator\_v5 and compute paired contrasts; (B) multi‑seed confirmatory design --- draw K >= 3 independent seeds per condition per task and estimate per‑condition pass probabilities with bootstrap CIs; (C) validator diversification --- add a second independent verifier (e.g., reachability emulator) and report intersection/union pass rates to reduce validator--task coupling. Because the generator, provider traces, and scoring harness are fully archived (execution/model\_configuration.json SHA‑256 a40e96e2..., execution/provider\_calls/*.json, execution/scoring/*.json), these reruns are deterministic and reproducible.

Threats to validity (artifact‑identified). Internal validity: the shared\_initial design intentionally removed generation variance and therefore precluded testing independent generative reliability. Construct validity: template‑aligned tasks and a single deterministic validator increase the chance that canonical candidates satisfy required\_sequences, possibly overstating success in open distributions. External validity: a single declared model (gpt‑5‑mini) and synthetic tasks limit generalization to other model families, open natural language intents, and production topologies. Conclusion validity: a degenerate contingency table produces trivial but correct inferential outputs; interpreting them beyond their precise estimand would be inappropriate. All these threats are documented in the published traces.

Operational implications for NetOps practice. Our artifact‑anchored diagnosis shows that (1) verifier grounding is necessary but not sufficient---experimental semantics must match the deployment question; (2) transparent preregistration plus exhaustive artifact publication materially clarifies what a given confirmatory test measures; and (3) simple rerun protocols (independent sampling, multi‑seed replication, validator diversity) convert a degenerate confirmatory design into an informative operational comparison without sacrificing reproducibility.

\section{Limitations}
\begin{itemize}
\item Controlled synthetic tasks: programmatic templates produced narrow required\_sequences and workflow\_policies that favour canonical solutions and limit external generalizability to heterogeneous operational intents.
\item Confirmatory shared\_initial semantics: the preregistered generation\_semantics = 'shared\_initial\_candidate' supplied identical canonical candidates to both conditions per pair, reducing independence between conditions and causing the degenerate paired outcome.
\item Single canonical seed for confirmatory test: the primary analysis used one canonical seed per task; preregistered exploratory multi‑seed analyses (seeds\_1..4) were not included in the confirmatory inference reported here.
\item Validator--task coupling: the deterministic validator enforces constraints encoded in the task templates, which can increase pass rates for template‑aligned candidates and reduce construct validity for open distributions.
\item Possible training‑data exposure: preregistered contamination heuristics found no flags under chosen thresholds, but absence of a flag is not proof of zero exposure to related artifacts in model pretraining.
\item Single‑model, single‑validator run: results apply only to gpt‑5‑mini under the recorded configuration and deterministic\_netops\_validator\_v5; they do not generalize across model families, agentic pipelines, or other validators.
\end{itemize}

\section{Conclusion}
We executed a fully preregistered, verifier‑grounded paired experiment and released a complete artifact bundle enabling byte‑level reproduction (execution/execution\_manifest.json and analysis/*). Under the preregistered confirmatory semantics (shared\_initial\_candidate) both the deterministic template baseline and the gpt‑5‑mini guarded condition validated on all N = 120 tasks (both 120/120). Archived evidence (shared\_initial candidate files and validator scoring traces) demonstrates that identical canonical candidates per pair and template‑tight task constraints caused the degenerate paired outcome; consequently the confirmatory paired result does not provide evidence about independent LLM generative reliability in operational settings. We provide artifact‑grounded, runnable modifications (independent\_per\_condition sampling, multi‑seed confirmatory execution, validator diversification, task heterogeneity expansion) that preserve reproducibility and enable informative independent comparisons. All execution and analysis artifacts---including the immutable master prompt reference (provenance/master\_prompt.txt; SHA‑256 1872df1e1805d2d96940456ca016bd665d1d5196add77f5acdf1582bb39b15ba) and the preregistration SHA‑256---are included in the submission bundle to permit exact audit and follow‑up experiments. Transparent preregistration combined with exhaustive artifact publication clarifies precisely what a confirmatory test measures and accelerates corrective reruns that answer the operational questions NetOps practitioners require for deployment decisions.

\section*{Disclosure Statement}
Mandatory disclosure (CNSM Agentic AI Researcher track). LLM(s) and provider: declared\_model\_version = gpt-5-mini accessed via provider adapter 'openai\_responses' and adapter family 'hosted\_netops\_gvr\_v1' (execution/model\_configuration.json SHA‑256 a40e96e212ab87da251ac0d6dbb75bc543afa13438b5a6b9ebd073597ffdd820). Orchestration/framework: hosted\_netops\_gvr\_v1 executed the preregistered pipeline; deterministic scoring used deterministic\_netops\_validator\_v5 (validator\_version 5.0). Preregistration: CAND-2026-001-NetConfig-Semantic-v1 (frozen prior to execution); preregistration SHA‑256 7db1663cf3982c8ea1e16c3dd7c0348f356bb4cbe3978ecc75bb59fdedb71a37. Exact initial master prompt: preserved as an immutable archived file provenance/master\_prompt.txt (SHA‑256 1872df1e1805d2d96940456ca016bd665d1d5196add77f5acdf1582bb39b15ba). Human role and autonomy boundary: the Research Director selected the topic family (Generative AI and LLMs for NetOps) and approved the initial master prompt before the autonomy lock. After the autonomy lock the autonomous pipeline performed literature review, research‑gap selection, preregistration, code generation, experiment execution, deterministic validation, statistical analysis, autonomous internal peer review, manuscript drafting and revision. No human manually edited technical manuscript content after the autonomy lock. Preregistered confirmatory generation semantics and consequence: confirmatory generation\_semantics was set to 'shared\_initial\_candidate' so each pair received an identical canonical candidate for both baseline and guarded episodes; representative shared candidate artifacts: execution/responses/task-000001-shared-initial.txt (SHA‑256 8f38c8becd7e1e36...), execution/responses/task-000002-shared-initial.txt (SHA‑256 ae967f70dfb6c...), execution/responses/task-000003-shared-initial.txt (SHA‑256 f7f4db249ecbbce...). Validator and scoring: deterministic\_netops\_validator\_v5 performed normalized parsing, intent constraint checking, transient safety checks, and emitted per‑episode JSON scoring traces archived under execution/scoring/*.json (representative SHA‑256s: execution/scoring/task-000001-guarded.json SHA‑256 0d56c1add7c5a57f...). Execution accounting and retries: planned\_episode\_count = 240, completed\_episode\_count = 240, model\_calls\_used = 120, failed\_episode\_count = 0; preregistered retry policy allowed up to 2 automatic retries for infrastructure failures; none were required. Contamination checks and reproducibility: preregistered string‑overlap heuristics found no confirmatory flags under chosen thresholds, but absence of a flag is not proof of zero exposure to related artifacts in model pretraining. All artifacts, seeds, code, and per‑file SHA‑256 hashes required for exact reproduction are released in the submission bundle (execution/execution\_manifest.json contains the exhaustive manifest). The initial master prompt is referenced immutably by path and SHA‑256 above.

\begin{thebibliography}{99}
\bibitem{ref1} https://doi.org/10.1145/3603269.3604866
\bibitem{ref2} https://doi.org/10.23919/cnsm62983.2024.10814407
\bibitem{ref3} https://doi.org/10.1002/nem.70029
\end{thebibliography}

\end{document}
